% Generated draft: requires mathematical and source review.
\documentclass[11pt]{article}
\usepackage[margin=25mm]{geometry}
\usepackage{fontspec}
\setmainfont{Latin Modern Roman}
\usepackage{amsmath,amssymb,mathtools,unicode-math}
\setmathfont{Latin Modern Math}
\usepackage{enumitem,needspace,xurl,hyperref}
\setlength{\parindent}{0pt}
\setlength{\parskip}{5pt}
\setlength{\emergencystretch}{3em}
\newcommand{\N}{\mathbb{N}}
\newcommand{\Z}{\mathbb{Z}}
\newcommand{\Q}{\mathbb{Q}}
\newcommand{\R}{\mathbb{R}}
\newcommand{\C}{\mathbb{C}}
\newcommand{\F}{\mathbb{F}}
\newcommand{\A}{\mathbb{A}}
\newcommand{\PP}{\mathbb{P}}
\newcommand{\E}{\mathbb{E}}
\newcommand{\eps}{\varepsilon}
\newcommand{\dd}{\,\mathrm{d}}
\begin{document}
\section*{UnsolvedMath ID: 30003385}
\textbf{Source problem code:} catalogue entry
\subsection{Definitions and mathematical statement}
The supplied record asks for a convergence theory for numerical approximations of viscosity solutions of Monge--Amp\`ere boundary problems. It does not specify a class of equations, boundary conditions, or numerical schemes sufficient to formulate a single quantified convergence theorem.

A schematic representation of the intended boundary problem is
\[
F(x,u(x),Du(x),D^2u(x))=0\quad\text{in }\Omega,
\qquad \mathcal B[u]=0\quad\text{on }\partial\Omega,
\]
where \(\Omega\subset\R^n\) denotes an as yet unspecified domain, \(F\) an as yet unspecified Monge--Amp\`ere type operator, and \(\mathcal B\) unspecified boundary data or a boundary constraint. This display is a template, not a specification of the admissible operators or boundary conditions. In particular, the elliptic branch and the admissible solution class have not been supplied.

For a specified continuous degenerate elliptic operator \(F\), the usual interior viscosity interpretation tests a continuous function \(u\) against \(C^2\) functions \(\phi\). With the sign convention that \(F\) is nonincreasing in its symmetric matrix argument, a local maximum of \(u-\phi\) at an interior point \(x\) requires
\[
F(x,u(x),D\phi(x),D^2\phi(x))\leq 0,
\]
and a local minimum requires the reverse inequality. A viscosity solution satisfies both requirements. Applying this definition to the intended Monge--Amp\`ere problems requires specifying the operator, its elliptic interpretation, any convexity or other admissibility restriction, and the treatment of boundary conditions; the record supplies none of these choices.

In the classical differentiable convex potential setting, the second boundary-value condition prescribes the gradient image,
\[
Du(\Omega)=\Omega^*,
\]
for a prescribed target domain \(\Omega^*\subset\R^n\). This explains the terminology in a standard setting only. The record does not specify source and target domains, their regularity or geometry, or the interpretation of this condition for nonsmooth potentials or more general transport costs.

\textbf{Requested research objective.} Establish hypotheses on more general Monge--Amp\`ere equations, boundary conditions, and numerical schemes under which the numerical approximations converge to viscosity solutions, including problems with the second boundary-value condition. For a discretization parameter \(h\to0\), the intended conclusion has the schematic form
\[
u_h\longrightarrow u,
\qquad u\text{ solves the specified boundary problem in the specified viscosity sense}.
\]
The admissible data, quantifiers over schemes, reconstruction of discrete approximations, convergence topology, and whether convergence concerns the whole family or subsequences remain unspecified. Thus this is a statement of a research objective with unresolved formulation requirements, not a complete convergence theorem.

The supplied dated assessment labels the problem partially solved: it reports a convergence framework for the quadratic-cost second boundary-value problem and several discretizations, while describing the corresponding theory for fully general equations and boundary data as remaining open. This is supplied curation, not an independently verified mathematical or literature conclusion.
\subsection{Short English statement}
Develop a convergence theory showing when numerical approximations of more general Monge--Amp\`ere boundary problems approach viscosity solutions, including the second boundary-value problem. The source leaves the equation class, boundary interpretation, numerical schemes, and convergence notion unspecified. Its supplied partially solved label and reported quadratic-cost progress require independent review.
\subsection{Sources}
\begin{itemize}
\item Supplied local UnsolvedMath registry record, permanent ID \(30003385\), titled \emph{Numerical Convergence for General Monge--Amp\`ere Boundary Problems}; collection: Oberwolfach Reports Open Problems; exact source problem code: \texttt{catalogue entry}. The record gives proposed year 2017 and names no individual proposer.
\item Supplied source citation: \emph{Applications of Optimal Transportation in the Natural Sciences (2018)}, Oberwolfach Reports, pp.~339--416, DOI \texttt{10.4171/owr/2017/7}. \url{https://doi.org/10.4171/owr/2017/7}. The supplied citation's 2018 date differs from the proposed year and report identifier referring to 2017; no resolution is supplied.
\item Supplied literature source: Froese Hamfeldt, \emph{Convergence framework for the second boundary value problem}. \url{https://doi.org/10.1137/18M1201913}.
\item Supplied literature source: \emph{Oberwolfach Report 7/2017}. \url{https://ems.press/content/serial-article-files/46673}.
\item Literature assessment date: \texttt{2026-08-21}. This is supplied dated curation, not independently verified. It reports coverage of the quadratic-cost second boundary-value problem and several discretizations, and labels the fully general extension open. The supplied background explicitly calls for independent verification of time-sensitive status and cited claims.
\item The registry labels the statement \texttt{corrected\_verified} and says that mathematical content was preserved while notation, typography, and sentence structure were repaired for publication. These are supplied editorial assertions, not independent verification of this draft or of the mathematical formulation.
\end{itemize}
\end{document}
